\documentclass[../../main.tex]{subfiles}
\graphicspath{{../../image/}{../../figure/}}
\begin{document}

\section{守恒律方程组的 L$^{1}$ 理论}\label{section:ch06-方程组L1理论}

前几节中基于 L$^{1}$ 与 L$^{\infty}$ 估计对标量守恒律的成功处理，自然引发如下期望：类似途径对守恒律方程组可能同样有效。不幸的是，情况似乎并非如此。为洞察其困难所在，考虑可对称化守恒律方程组的 Cauchy 问题：
\begin{align}
\partial_{t}U+\sum_{\alpha=1}^{m}\partial_{\alpha}G_{\alpha}(U)=0,\qquad x\in\mathbb{R}^{m},\ t>0,
\label{align:ch06-6-10-1}
\end{align}
\begin{align}
U(x,0)=U_{0}(x),\qquad x\in\mathbb{R}^{m}.
\label{align:ch06-6-10-2}
\end{align}

对标量情形类比\cref{definition:ch06-6-2-1}，我们要求\eqref{align:ch06-6-10-1}、\eqref{align:ch06-6-10-2}的可容许解对任一凸熵--熵流对 $(\eta,Q)$（$\eta$ 凸）满足\eqref{align:ch04-熵积分0}。

第一个检验应当是：\eqref{align:ch06-6-10-1}的平凡常值解 $\bar{U}$ 在可容许解类中是否 L$^{p}$ 稳定：
\begin{align}
\|U(\cdot,t)-\bar{U}\|_{L^{p}(\mathscr{B}_{r})}\leqslant c_{p}\|U_{0}(\cdot)-\bar{U}\|_{L^{p}(\mathscr{B}_{r+st})}.
\label{align:ch06-6-10-3}
\end{align}
由于方程组可对称化，从而配备二次增长的凸熵，由\cref{theorem:ch05-involution-existence}，至少对 $p=2$，\eqref{align:ch06-6-10-3}成立。

于是问题产生：此类估计是否也可能对 $p\neq 2$ 成立，其中 $p=1$ 与 $p=\infty$ 情形特别引人关注。

对线性组
\begin{align}
\partial_{t}V+\sum_{\alpha=1}^{m}\mathrm{D}G_{\alpha}(\bar{U})\partial_{\alpha}V=0,
\label{align:ch06-6-10-4}
\end{align}
上述估计在任何 L$^{p}$ 中都成立，因为相应流形是沿特征线的平移。但是，非线性效应是否会保持这种稳定性，取决于线性化矩阵族 $\{\mathrm{D}G_{\alpha}(\bar{U})-zI\}$ 的谱，其中 $z\in\mathbb{R}$ 是参数。下面命题给出 $L^{p}$ 稳定性的完整刻画。

\begin{theorem}\label{theorem:ch06-6-10-1}
常值解 $\bar{U}$ 在可容许解类中是 L$^{p}$ 稳定的，当且仅当存在 $\bar{U}$ 的邻域 $\mathscr{B}_{\delta}(\bar{U})$，使得对每个 $z\in\mathbb{R}$，矩阵族 $\{\mathrm{D}G_{\alpha}(U)-zI\}$（$U\in\mathscr{B}_{\delta}(\bar{U})$）可通过与 $p$ 相关的实对角化同时对称化。
\end{theorem}

特别地，该条件对 $p=1$ 与 $p=\infty$ 只在线性组时成立。于是标量守恒律所享有的极强 L$^{1}$ 压缩性质，一般不能推广到方程组。

\begin{theorem}\label{theorem:ch06-6-10-2}
设 $m\geqslant 2$，$n\geqslant 2$。则存在 $\bar{U}\in\mathbb{R}^{n}$ 及 $\bar{U}$ 的邻域 $\mathscr{B}_{\delta}(\bar{U})$，使得对任一初值 $U_{0}\in\mathscr{B}_{\delta}(\bar{U})$，$\|U_{0}\|_{L^{\infty}(\mathbb{R}^{m})}\leqslant\delta$，\eqref{align:ch06-6-10-1}、\eqref{align:ch06-6-10-2}的可容许弱解 $U$ 在 $\mathbb{R}^{m}\times(0,\infty)$ 上是无界的。
\end{theorem}

其证明见\cref{section:ch06-注记}所引文献。


尽管存在上述障碍，仍有一类方程组满足\eqref{align:ch06-6-10-3}，对 $1\leqslant p\leqslant 2$，通过求解 $\mathscr{B}_{\delta}(\bar{U})$ 上的 Goursat 问题即可构造具有所需性质的熵。事实上，在对组附加额外假设下，甚至可能构造对所有 $p$ 都满足\eqref{align:ch06-6-10-3}的凸熵；对此类组，常值解在整个范围 $1\leqslant p\leqslant\infty$ 上是 L$^{p}$ 稳定的。

满足\eqref{align:ch06-6-10-3}的方程组类特别包括标量守恒律（$n=1$）（任意空间维数 $m$），以及单空间维数（$m=1$）中任意规模 $n$ 的方程组；但除此之外，它含有的（即使是适度的）物理意义上的代表非常少。一个例子是方程组
\begin{align}
\partial_{t}U+\sum_{\alpha=1}^{m}\partial_{\alpha}[F_{\alpha}(|U|)U]=0,
\label{align:ch06-6-10-9}
\end{align}
它描述流体在各向异性多孔介质中的流动。此方程组的特殊性质使其在分析上易于处理，因而可作为展示研究若干空间维数中双曲守恒律方程组所面临问题的载体。

若 $U$ 是\eqref{align:ch06-6-10-9}的经典解，则易见其“密度” $\rho=|U|$ 满足标量守恒律
\begin{align}
\partial_{t}\rho+\sum_{\alpha=1}^{m}\partial_{\alpha}[\rho F_{\alpha}(\rho)]=0,
\label{align:ch06-6-10-10}
\end{align}
而其方向单位向量场 $\Theta=\rho^{-1}U$ 满足输运方程
\begin{align}
\partial_{t}\Theta+\sum_{\alpha=1}^{m}F_{\alpha}(\rho)\partial_{\alpha}\Theta=0.
\label{align:ch06-6-10-11}
\end{align}
于是 Cauchy 问题\eqref{align:ch06-6-10-9}、\eqref{align:ch06-6-10-2}的经典解可这样构造：先以初值 $\rho(\cdot,0)=|U_{0}(\cdot)|$ 求解\eqref{align:ch06-6-10-10}（例如用\cref{section:ch06-cauchy问题经典解}所阐述的特征线法），再由 $\Theta$ 沿常微分方程轨迹保持常值这一性质来确定 $\Theta$：

\begin{align}
\frac{\mathrm{d}x}{\mathrm{d}t}=F(\rho(x,t)).
\label{align:ch06-6-10-12}
\end{align}

如何把上述程序推广到弱解并不显然。当然仍可把 $\rho$ 确定为\eqref{align:ch06-6-10-10}以初值 $|U_{0}|\in L^{\infty}$ 的可容许弱解，但当 $F(\rho(x,t))$ 仅是 $L^{\infty}$ 函数时应如何解释\eqref{align:ch06-6-10-12}，绝非清楚。事实上，已证明（参见\cref{section:ch06-注记}）\eqref{align:ch06-6-10-9}的 Cauchy 问题一般在 $L^{\infty}$ 中不适定。

一个相关而有力的常微分方程理论 $\dot{X}=P(X)$ 存在，但它要求 $P$ 是 $BV$ 中散度为零的向量场。为应用该理论，限制初值使 $|U_{0}|$ 是 $\mathbb{R}^{m}$ 上局部有界变差的正函数。由\cref{theorem:ch06-6-2-3}与\cref{theorem:ch06-6-2-6}，这将保证 $\rho$ 在上半空间上是局部有界变差的正函数。接着重标度时间变量，把\eqref{align:ch06-6-10-12}改写为隐式形式
\begin{align}
\begin{cases}
\dfrac{\mathrm{d}t}{\mathrm{d}\tau}=\rho(x,t),\\[6pt]
\dfrac{\mathrm{d}x}{\mathrm{d}\tau}=\rho(x,t)F(\rho(x,t)),
\end{cases}
\label{align:ch06-6-10-13}
\end{align}
它具有所期望的特征：由\eqref{align:ch06-6-10-10}，向量场 $(\rho,\rho F(\rho))$ 在上半空间上散度为零。

通过在\eqref{align:ch06-6-10-13}的解族 $(t(\tau),x(\tau))$ 中消去 $\tau$，得到曲线族 $x=x(t)$，即\eqref{align:ch06-6-10-12}的形式轨迹，沿其 $\Theta$ 保持常值。因而 $\Theta$ 可由其初值确定，而初值可仅属于 $L^{\infty}$。最后可证（文献见\cref{section:ch06-注记}）$U=\rho\Theta$ 是\eqref{align:ch06-6-10-9}、\eqref{align:ch06-6-10-2}的弱解：


\begin{theorem}\label{theorem:ch06-6-10-3}
对任一 $U_{0}\in L^{\infty}(\mathbb{R}^{m};\mathbb{R}^{n})$ 且 $|U_{0}|\in BV_{\mathrm{loc}}(\mathbb{R}^{m})$，在 $[0,\infty)$ 上存在\eqref{align:ch06-6-10-9}、\eqref{align:ch06-6-10-2}的唯一 $L^{\infty}$ 弱解 $U$，使 $\rho=|U|$ 是\eqref{align:ch06-6-10-10}以初值 $\rho(\cdot,0)=|U_{0}(\cdot)|$ 的可容许弱解。
\end{theorem}

进而已证明，上述解在 $L^{1}_{\mathrm{loc}}$ 中连续依赖于其初值，并且对\eqref{align:ch06-6-10-9}的任一凸熵满足熵可容许性条件，至少当函数 $\rho F(\rho)$ 在 $(0,\infty)$ 上的临界点集测度为零时成立。

另一方面，即使初值具有小总变差，\eqref{align:ch06-6-10-9}的 Cauchy 问题在 $BV$ 中也是不适定的：

\begin{theorem}\label{theorem:ch06-6-10-4}
设 $m\geqslant 3$，$n\geqslant 2$。对任一非零 $\bar{U}\in\mathbb{R}^{n}$（使 $|\bar{U}|$ 不是 $\rho F(\rho)$ 的临界点）以及任意 $\delta>0$，存在初值 $U_{0}$，取值于 $\mathscr{B}_{\delta}(\bar{U})$，当 $|x|>1$ 时等于 $\bar{U}$，在 $\mathbb{R}^{m}$ 上总变差小于 $\delta$，并具有如下性质：若 $U$ 是\eqref{align:ch06-6-10-9}、\eqref{align:ch06-6-10-2}在某时间区间 $[0,T)$ 上的任一可容许 $L^{\infty}$ 弱解，则 $U$ 在 $\mathbb{R}^{m}\times[0,T)$ 上的总变差为无穷。
\end{theorem}

读者应牢记，\eqref{align:ch06-6-10-9}如此特殊，以至于上述结论不必解释为一般方程组行为的代表。若干空间变量中双曲守恒律方程组的理论仍处于初期阶段。

\end{document}
